\section{Introdução}

O 1º Hackathon do curso de tecnologia do IFPR Campus Pinhais foi organizado por uma comissão pequena e voluntária, formada pela docente de Engenharia de Software I, um monitor da disciplina e dois representantes de um projeto de extensão. Antes do sistema, as inscrições eram feitas por planilha e e-mail, o que já havia gerado inscrições duplicadas, a avaliação dos jurados era em papel e a apuração do resultado levava horas e era sujeita a contestação.

Como projeto da disciplina, a turma construiu um sistema web único que cobre o ciclo completo do evento: inscrição, formação de equipes, agenda, \textit{check-in} por QR, submissão de projetos, avaliação por jurados, resultados, voto popular e certificados em PDF. O código está em um repositório único, em produção na plataforma Railway, e foi desenvolvido entre 6 de agosto e 31 de outubro de 2026.

Os autores participaram do desenvolvimento do sistema. Esta é, portanto, uma análise interna, e não de um projeto de terceiros. Essa condição permite acesso completo às decisões e às evidências, mas exige cautela com o viés, que tratamos na Seção~\ref{sec:met}. A escolha do projeto foi acordada com a docente da disciplina.

O objetivo do trabalho é analisar criticamente o projeto sob a ótica da Engenharia de Software, respondendo a quatro questões de pesquisa:

\begin{itemize}
  \item \textbf{QP1:} Que processo de desenvolvimento o projeto adotou e como ele se comportou na prática?
  \item \textbf{QP2:} Como a gerência de configuração e de mudanças foi conduzida?
  \item \textbf{QP3:} Que evidências existem de qualidade de produto e de processo?
  \item \textbf{QP4:} Como o projeto organizou testes e manutenção?
\end{itemize}

O restante do artigo está organizado da seguinte forma. A Seção~\ref{sec:fund} apresenta a fundamentação teórica. A Seção~\ref{sec:met} descreve a metodologia. A Seção~\ref{sec:res} traz os resultados. A Seção~\ref{sec:disc} discute pontos fortes, fragilidades e recomendações, e a Seção~\ref{sec:conc} conclui.
